\begin{statementbox}
\textbf{C004\quad 椭圆顶角最大值（长轴两端点的张角）}

设非圆椭圆
\[
C:\quad \frac{x^2}{a^2}+\frac{y^2}{b^2}=1,\qquad a>b>0,
\qquad c^2=a^2-b^2,
\]
其长轴两端点为 $A_1(-a,0)$、$A_2(a,0)$。取
$P(x_0,y_0)\in C\setminus\{A_1,A_2\}$，令
$\alpha=\angle A_1PA_2$（\textbf{普通内角}，$0<\alpha<180^\circ$）。

\begin{enumerate}
\item \textbf{角的范围、最大值与位置（核心结论）}：
\[
\boxed{90^\circ<\alpha\le\alpha_{\max}
=2\arctan\frac{a}{b}<180^\circ.}
\]
当且仅当 $P=(0,b)$ 或 $(0,-b)$（\textbf{短轴端点}）时，
取到最大值。等价地
\[
\boxed{\alpha_{\max}=90^\circ+
\arctan\frac{a^2-b^2}{2ab}},\qquad
\boxed{\tan\alpha_{\max}=-\frac{2ab}{a^2-b^2}}.
\]
注意 $\alpha_{\max}$ 是\textbf{钝角}，故其正切值为\textbf{负}。

\item \textbf{最小值（严格区分“最小值”与“下确界”）}：
在上述定义域中，$\alpha$ \textbf{没有最小值}；
当 $P$ 沿椭圆趋近长轴两端点时，$\alpha\to90^\circ$，
故下确界为 $90^\circ$。在端点处 $\angle A_1PA_2$ 无定义，
不可说“在长轴端点取得最小值”。

\item \textbf{任意动点的角度--高度关系（最推荐的通式）}：
\[
\boxed{\cot\alpha=-\frac{(a^2-b^2)|y_0|}{2ab^2}},
\qquad
\boxed{\alpha=90^\circ+
\arctan\frac{(a^2-b^2)|y_0|}{2ab^2}}.
\]
在 $a,b$ 固定时，$\alpha$ 随 $|y_0|$ \textbf{严格增大}。
特别地（$y_0\ne0$），
\[
\boxed{\tan\alpha=-\frac{2ab^2}{(a^2-b^2)|y_0|}}.
\]

\item \textbf{斜率乘积定值}：
\[
\boxed{k_{PA_1}\,k_{PA_2}=-\frac{b^2}{a^2}},\qquad
k_{PA_1}=\frac{y_0}{x_0+a},\quad
k_{PA_2}=\frac{y_0}{x_0-a}.
\]
由于已排除两个长轴端点，上述斜率均有定义。

\item \textbf{等角、等面积、等高}：
设 $S=S_{\triangle A_1PA_2}$，则
\[
\boxed{S=a|y_0|=\frac{2a^2b^2}{a^2-b^2}
\tan(\alpha-90^\circ)}.
\]
在\textbf{同一椭圆}上，$\alpha$ 相同 $\Longleftrightarrow$
$S$ 相同 $\Longleftrightarrow$ $|y_0|$ 相同，并且 $|x_0|$ 也确定。
通常对应四个轴对称点；当 $\alpha=\alpha_{\max}$ 时只有
两个短轴端点。
\end{enumerate}

\textbf{圆的边界情形：}若 $a=b$，椭圆退化成圆，
由直径所对圆周角得 $\alpha\equiv90^\circ$；
此时 $\tan\alpha$ 不存在，上述含 $a^2-b^2$ 作分母的公式不可套用。
\end{statementbox}
